\majorbreak
\section{Simplify the interfaces between proofs}
\label{sec:architecture}

The strongest abstractions here follow repeated mathematical boundaries:
coordinate packing, source-window extraction, and degree-layer estimates.
They should reduce the number of arguments a maintainer has to understand,
not merely hide an unchanged long proof behind a new name. The proposed
helpers in this section are not existing declarations.

\subsection{Prove the nine coordinate packaging equivalence once}
\label{sec:nine}

\textbf{Priority 1.} The blocks in
\src{SplitSourceExistence.lean}{98}{157} and
\src{GapSourceExistence66.lean}{53}{112} are byte-identical: each is 60 lines
and 2,292 UTF-8 bytes, excluding a trailing newline. They define positive and
natural coordinate interpretations and prove the two existence equivalences.
The scalar evaluator comes from the surrounding namespace; the packing
argument does not inspect its arithmetic.

\existing{} Gap already reuses Split's cofinal-cap lemmas at
\src{GapSourceExistence66.lean}{12}{15}. That is a successful abstraction to
preserve. The remaining duplicate is specifically the coordinate packaging.

\proposed{} Parameterize one helper by the evaluator below, with a predicate
\code{HasZero} expressing the same positive-coordinate existential. Retain
the old public definitions and theorem types; specialize the generic proof
inside their bodies.

\begin{Lean}
-- Proposed helper boundary, UNCOMPILED.
abbrev Evaluator :=
  ℕ → ℕ → FreeCarrierRatio.FreeVars → ℕ → ℤ

def positiveVars (x : Fin 9 → ℕ) : FreeCarrierRatio.FreeVars :=
  ⟨x 2, x 3, x 4, x 5, x 6, x 7 - 1⟩

def naturalVars (x : Fin 9 → ℕ) : FreeCarrierRatio.FreeVars :=
  ⟨x 2 + 1, x 3 + 1, x 4 + 1, x 5 + 1, x 6 + 1, x 7⟩
\end{Lean}

This captures an important asymmetry. The positive tuple stores
\code{b,g,h,s,w,C,i,j+1,n}, whereas the internal \code{FreeVars.j} is an
unrestricted natural. In the natural tuple, slot 7 is \code{j}, not
\code{j+1}. The generic theorem must not add \code{0 < v.j}, or omit
positivity of \code{g} or \code{n}.

The supplied \pathcode{architecture/NineCoordinatePackaging.lean} contains
the complete proposed helper, including the successor/predecessor bridge.
The companion compatibility file states four examples matching the two
existing positive/natural existence interfaces in each source package.
A conceptual replacement is:

\begin{Lean}
-- Schematic specialization; helper name is proposed.
exact NineCoordinatePackaging.exists_natural_zero_iff_hasZero
  (GapSourceExistence66.value d a)
\end{Lean}

\risk{} This is a good first abstraction because it moves no number theory
and assumes no source correctness. Put it below both packages with only the
needed foundational imports. Compile the helper and all four compatibility
examples first. The 60-line duplicate is an observed opportunity, not a net
60-line saving: the generic interface and compatibility wrappers also cost
source text. The main gain is one maintained coordinate-bijection proof.

\subsection{Give long conjunctions names without changing their definitions}
\label{sec:accessors}

\textbf{Priority 1.} \code{Carrier66.Canonical} is a seven-part conjunction
at \src{Carrier66.lean}{153}{163}. A consumer in
\code{canonical_strong_error} reaches a fact through six positional
projections at \src{Carrier66Source.lean}{167}{191}. The fragility is not
proof search; it is that a new conjunct can silently make every positional
access harder to read or maintain.

\begin{Lean}
-- Current source, Carrier66Source.lean:180.
rw [hc.2.2.2.2.2.1, sub_self] at hid

-- Proposed accessor, after proving Canonical.Z_eq_delta.
rw [hc.Z_eq_delta, sub_self] at hid
\end{Lean}

Add derived projection lemmas such as \code{Canonical.Z_eq_delta},
\code{Canonical.index_eq_one}, and \code{Canonical.rank_pos}. Initially
keep \code{Canonical} itself unchanged, so existing destructuring patterns
and constructors continue to work. A later \code{Prop} structure is optional,
with a proved equivalence; it is not required to obtain readable accesses.

The same restrained approach can expose \code{PortBounds.beta_pos} and a
small core-positivity interface. Reuse \code{Carrier65.scales_pos} at
\src{Carrier65.lean}{175}{183} rather than turning derived positivity into
new assumptions. The existing \code{CurrentPolynomialWindow2011.Facts}
record at \src{CurrentPolynomialWindow2011.lean}{27}{68} is the local
precedent for coherent named fields.

\subsection{Share one coherent source and scalar window certificate}
\label{sec:window}

\textbf{Priority 2.} The low-digit and high-digit consumers repeat the same
restoration and normalization pipeline. Compare
\code{polynomial_zero_reconstructs} at
\src{CurrentSparseDigits2011.lean}{111}{138} with
\code{polynomial_zero_mask_test} at
\src{CurrentHighDigits2011.lean}{135}{159}. Their 13-line preambles turn a
polynomial zero into a \code{Carrier65.value} zero, restore a canonical
parent, call the scalar-window theorem, rewrite the restored core, and
normalize powers into the current radix. The fuller
\code{polynomial_zero_window} repeats the parent/source bridge because it
also needs the actual source assignment
(\src{CurrentPolynomialWindow2011.lean}{78}{119}).

\existing{} \code{Carrier66Source.canonical_window} at
\src{Carrier66Source.lean}{376}{394} already proves the generic scalar
window. \code{positive_carrier_zero_source} at
\src{Carrier66Source.lean}{398}{416} already supplies a source lift.
The proposal is to compose and expose these results once, not to replace
them with a second theory of windows.

\proposed{} First expose a direct positive \code{Carrier65} zero-to-source
certificate containing the natural source, positivity, rows, preserved outer
coordinates, and the existing \code{sourceQ_eq}, \code{gap_eq}, and
\code{nativeRank_eq} equalities. Then add one current-mask specialization
that normalizes width/radix and presents both mask tests together.

\begin{Lean}
-- Interface sketch only; names and record are proposed.
obtain ⟨N, hN, facts⟩ := current_zero_window_facts K hK v hv hz
-- Low decoder consumes facts.C_pos, facts.C_split, facts.low_mask.
-- High decoder consumes facts.high_mask.
-- Both use the same facts.q_window and facts.J_window.
\end{Lean}

The record must carry one \emph{shared} \code{N} and source assignment.
Separate existential choices for different fields would weaken the usable
interface. Keep source lift, scalar window, and later carry/digit facts in
distinct layers so the low-level helper does not import its consumers.

\risk{} Preserve the high port's single beta shift and both origin
corrections, low \code{+1} and high \code{-1}. The public theorem must still
derive the certificate from positivity and the actual polynomial zero.
Adding a canonical-parent, source-zero, no-carry, or decoded-window
hypothesis would shorten the proof by changing its theorem, not simplify it.

\subsection{Put the current base interface below completeness}

\textbf{Priority 3, after the window interface.}
\src{CurrentCarrierArithmetic2011.lean}{1}{2} imports both
\code{CurrentMasks2011} and \code{ThirteenCarrierComplete}. Its low/high
digit clients need the four-port bounds, \code{polynomial65}, and the
carrier-to-source arithmetic bridge, not the thirteen-coordinate completeness
layer. The latter's definitions and equivalences occur at
\src{CurrentCarrierArithmetic2011.lean}{45}{132}.

Extract a proposed \code{CurrentCarrierBase2011} below completeness,
importing \code{CurrentMasks2011} and \code{Carrier66Source}. Keep the old
arithmetic module as a compatibility facade and redirect only the simpler
clients to the base. The static local import closure changes from 37 existing
modules to 31 existing modules plus the proposed base. The five avoided
existing dependencies, apart from replacing the old facade itself, are
\code{SignedPellResidue}, \code{ThirteenCarrier},
\code{ThirteenCarrierComplete}, \code{ThirteenCarrierSoundness}, and
\code{ThirteenFreshSeed}.

This is a measured graph comparison, \emph{not} a build-time improvement.
External transitive imports were not traversed. Preserve the facade's names
and avoid a dependency cycle by never importing the digit/window consumers
back into the base.

\subsection{Abstract the degree profile but retain sharp hypotheses}
\label{sec:profiles}

\textbf{Priority 3.} The source-positivity arguments in
\src{SplitWordBounds.lean}{76}{123} and
\src{GapWordBounds66.lean}{31}{78} share the same numerical proof structure.
Their upper-scale calculations likewise appear at
\src{SplitWordBounds.lean}{149}{200} and
\src{GapWordBounds66.lean}{81}{132}. Gap already aliases several source and
packing definitions and imports common geometric lemmas at
\src{GapWordBounds66.lean}{11}{20}; preserve that reuse.

The proposed generalization is a small degree profile $(d_0,d_1,d_2)$ with
the two actual gap inequalities
\[
 d_1 < d_2 + 116,\qquad d_0 < d_2 + 232.
\]
Retain the leading-coefficient condition \code{D2.coeff d2 = 1}, each
degree and height bound, the scale $B=4096Hb^2$, and $B^{116}\le U$.
Prove lower-layer domination and positivity once. Instantiate the profile
with $(1719,2027,1929)$ for Split and $(1627,1941,1850)$ for Gap.

\begin{Lean}
-- Schematic proposed API, not an existing theorem.
exact degreeProfile_negative_layers hk.profile hH hb hscale hU
  (by decide : 1941 < 1850 + 116)
  (by decide : 1627 < 1850 + 232)
\end{Lean}

Only after this lower-bound interface proves useful should it absorb the
upper packing estimate. A target exponent $d_2+233$ gives the existing
concrete exponents 2162 and 2083. Do not replace either sharp bound by a
larger one simply to make the abstraction easier.

\risk{} Gap's \code{KernelBounds} is not a stronger version of Split's:
its leading coefficient is at 1850 rather than 1929. Compare
\src{SplitWordBounds.lean}{20}{27} and
\src{GapWordBounds66.lean}{22}{29}. An attempted coercion between them
would be mathematically wrong. Store degree data and elementary inequalities,
not the final conclusions. Symbolic exponents can also increase elaboration
cost, so benchmark both concrete specializations, not only the new theorem.

\subsection{Remove duplicated coefficient expressions in stages}
\label{sec:coefficients}

\textbf{Priority 4.} The region
\src{NativeNinePolynomial52.lean}{93}{404} contains 22
\code{mapRowCoefficient*} helpers and their dispatch, with repeated
\code{hleft}/\code{hright} spellings of original and mapped coefficients.
The authoritative table is only
\src{NativeQuadraticParent.lean}{155}{177}. Gap has analogous transport
cases at \src{GapNinePolynomial66.lean}{99}{283}.

The reviewed native table uses fixed port indices, integer numerals,
negation, ring addition, subtraction and multiplication, and the fixed natural
exponent 2. Its branches compare a fixed pair index, not ring values. That
specific fragment admits a ring-homomorphism interpretation. It does not
justify transporting arbitrary natural evaluators, truncated subtraction,
division, or \code{toNat}. The natural exponent subtraction in
\code{kernelValue} at \src{NativeQuadraticParent.lean}{64}{64} is fixed
layout metadata and must remain fixed.

\existing{} Generic evaluator naturality already exists in
\src{ArithmeticDAG.lean}{48}{103}; the carrier and polynomial files already
use \code{value_map} and \code{eval_polynomial}, for example
\src{Carrier65.lean}{155}{167} and
\src{NativeNinePolynomial52.lean}{406}{454}.

\proposed{} Try the cheaper change first: a focused local simp interface
or generated case tactic that keeps the coefficient table and removes the
duplicated intermediate expressions. Only if measured maintenance savings
warrant it, encode these small coefficients once as computable ring syntax
and prove its evaluator naturality once. Reuse the existing DAG evaluator
where it helps; do not force a tiny coefficient into a large framework.

\risk{} Keep a proved equality to the old authoritative table, preserve
the general \code{CommRing} theorem, and retain row order, support,
repeated-term accumulation, and corrections. Do not define both the source
and certificate from the same replacement table and treat reflexivity as
independent verification. Replacing the ordinary integer evaluator with
noncomputable polynomial evaluation would cross the intentional boundary
documented at \src{NativeNinePolynomial52.lean}{9}{13}.
